\documentclass[11pt]{article}
\usepackage[margin=.88in]{geometry}
\usepackage{amsmath,amssymb,amsthm,booktabs,hyperref,microtype}
\hypersetup{colorlinks=true,linkcolor=black,urlcolor=blue,citecolor=black}
\newtheorem{theorem}{Theorem}[section]\newtheorem{lemma}[theorem]{Lemma}\newtheorem{corollary}[theorem]{Corollary}
\newcommand{\sig}{\operatorname{sig}}
\title{Acyclic matchings and sharp line-graph signatures beyond nine cycles}
\author{Blake Stone's mathematics project\\\small AI-assisted research manuscript}
\date{November 27, 2026}
\begin{document}\maketitle
\begin{abstract}
Let $f(c)$ be the largest adjacency signature of the line graph of a connected simple graph with $c$ independent cycles. We prove $f(10)=5$, $f(11)=6$, $f(12)=6$ and $f(13)=7$ and prove $f(c)\le\lfloor2c/3\rfloor$ for every $c\ge2$, so $\lim_{c\to\infty}f(c)/c\le2/3$. The previous upper bound was $5/6$. The Paone--Paone conjectured upper bound holds at every cyclomatic number for simple branch kernels, with arbitrary subdivisions and attached trees. Our obstruction is an acyclic matching in a cubic expansion. On a suitable subspace the quadratic form becomes the adjacency form of a forest even after signed identifications. Exact finite computations treat the remaining small-cycle cases.
\end{abstract}
\section{Statements and the cubic form}
Write $s(L(G))$ for the number of positive adjacency eigenvalues minus the number of negative ones, counted with multiplicity, of the line graph of a finite connected simple graph $G$. Its cyclomatic number is $c=|E(G)|-|V(G)|+1$. Paone and Paone conjectured that the largest possible signature at fixed $c\ge1$ is $f(c)=\lfloor(c+1)/2\rfloor$ and gave matching lower constructions \cite[Theorem 5.3 and Conjecture 5.4]{PP}. The earlier manuscripts \cite{prev,finite} established the formula through $c=9$ and showed that $\alpha=\lim f(c)/c$ exists, between $1/2$ and $5/6$.
\begin{theorem}\label{main}
For every $c\ge2$ one has
\[
 f(c)\le \lfloor 2c/3\rfloor,
 \qquad \tfrac12\le\alpha\le\tfrac23.
\]
For a root graph with a simple branch kernel and $c\ge5$, $s(L(G))\le\lfloor(c-1)/2\rfloor$. At all cyclomatic numbers the source's upper bound holds for such kernels. It also holds whenever the compressed underlying graph of a cubic expansion, defined in Section~\ref{compression}, is a cycle, or a tree with a matching that misses at most one vertex. The bound is sharp among path compressions at every $c\ge4$. For $c\ge9$, an even-rank counterexample would need at least six loop vertices and double-pair units in total in every cubic expansion; an odd-rank counterexample would need at least eight.
\end{theorem}
\begin{theorem}\label{finite-theorem}
For every $c\ge10$ one has $f(c)\le c-5$, and for every $c\ge12$ one has $f(c)\le c-6$. In particular $f(10)=5$, $f(11)=6$, $f(12)=6$, and $f(13)=7$. The next case satisfies $7\le f(14)\le8$.
\end{theorem}
To recall the finite reduction \cite[Sections 2--3]{prev}, remove pendant trees and suppress degree-two paths to obtain the branch kernel $K$ of the $2$-core. For $c\ge2$ this is a connected multigraph of minimum degree at least three. Split a high-degree vertex of degree $d$ into two adjacent vertices distributing two and $d-2$ of the old incidences between them. Iteration produces a connected cubic multigraph $C$ of order $J=2c-2$, called a cubic expansion; loops count twice at a vertex. The resulting relaxation bounds the quantity $\mu=s(L(G))+c-1$ as follows.
\begin{lemma}\label{form}
At each vertex of $C$ start with diagonal coefficient one. Each edge $uv$ gives a term $2\epsilon x_ux_v$, an equation $x_u=\epsilon x_v$, or deletion at cost one, with independent signs $\epsilon=\pm1$. Each paid charge at a vertex adds two to its diagonal at cost one. Then $\mu$ is at most the maximum signature of the form on the equation space, less all costs. A further upper bound permits arbitrary coordinate pins $x_u=0$ and uses only pins, equations and signed terms.
\end{lemma}
\begin{proof}
The branch-vertex relaxation of \cite[Lemma 3.1]{prev} gives the first assertion on $K$ with initial diagonals $d_K(u)-2$. It includes arbitrary pendant trees and all subdivisions. Replace a branch of degree $d$ by its cubic splitting tree. It has $d-2$ vertices, since if it has $k$ then $3k=d+2(k-1)$. Set every splitting edge to the equation of equality. The cubic diagonals now sum to the original diagonal and all other choices pull back. This proves the first assertion on any expansion.

Restricting any symmetric form to a subspace of codimension at most one can decrease its signature by at most one, by interlacing. For a paid charge, pin its coordinate and remove the charge. For a deletion, pin one of its ends, then restore the edge as an arbitrary signed term. The added pin has codimension at most one on the current equation space and the removed cost compensates the possible signature loss. Applying this to each payment proves the second assertion.
\end{proof}
\section{Acyclic matching obstruction}
In a multigraph a qualifying matching here consists of disjoint single nonloop edges joining vertices of unit diagonal without loops. Call it acyclic when the entire multigraph induced by its endpoints is a forest. In particular parallel edges between selected endpoints are excluded.
\begin{lemma}\label{forest}
An acyclic qualifying matching of size $h$ in a cubic expansion implies $s(L(G))\le c-1-h$. More generally, in the pin-and-equation relaxation on $N$ coordinates with arbitrary initial diagonal coefficients, such a matching on unit diagonal vertices bounds the final signature by $N-h$.
\end{lemma}
\begin{proof}
Fix the signs, equations and pins of Lemma~\ref{form} and write $W$ for their solution space. For each matched odd edge $uv$ choose the line $x_u=-\epsilon x_v$ on its two coordinates. The two diagonal squares and its signed edge term then sum to zero. For a matched edge giving an equation $x_u=\epsilon x_v$, use its opposite line $x_u=-\epsilon x_v$; the coordinates vanish on intersection with $W$. The direct sum of these $h$ lines, with all other coordinates zero, is a subspace $V$ of dimension $h$.

Contract the matched pairs in the forest induced by their endpoints. The resulting graph on at most $h$ variables is still a forest. All terms left on $W\cap V$ come from its other internal odd edges. The other internal equations identify their two variables up to sign, while an equation to a coordinate outside $V$ or a pin sets the incident variable to zero. Contracting the equation edges, absorbing signs in the free variables, and deleting the zero variables leaves a forest. There is no diagonal term after these operations: a surviving odd edge inside an identified class would give a cycle before contraction. The restricted form has signature zero, since negating all variables of one part of a bipartition takes it to its negative. Finally $\dim W-\dim(W\cap V)\le N-h$. Interlacing gives $\sig(F|_W)\le N-h$. On $C$, $N=2c-2$, proving the claim.
\end{proof}
As a sharpness check on the relaxation, a forest $T$ with all unit diagonals and arbitrary free pins has maximum signature $|V(T)|-\nu(T)$, where $\nu$ is the ordinary matching number. Every matching is acyclic, giving the upper bound. Pin a minimum vertex cover, of size $\nu(T)$ by bipartite matching duality. All remaining variables have independent positive squares, attaining the bound.

\section{Compression and large graphs}\label{compression}
For $J\ge6$, delete the cubic vertices that carry a loop, each of which has a single link to the remainder. Contract each double pair to a single vertex and discard its two internal edges. Triple edges are impossible in a connected cubic graph on more than two vertices. The resulting connected loopless subcubic multigraph is $S$, with underlying simple graph $U$. Write $n=|V(S)|$, $\ell=|E(S)|-n+1$, and $e=|E(U)|-n+1\le\ell$. As in \cite[Section 2]{prev}, each loop-vertex reduction and each double-pair reduction decreases both the order and the cyclomatic number by one; hence $n=c-2+\ell$.

Any acyclic matching of $U$ lifts to one in $C$. For each matched adjacency choose an original single link between the two units. At most one original vertex is selected from each unit. Between selected vertices every surviving adjacency maps to one in the induced forest of $U$; an original double edge is internal to its unit and an original loop vertex was removed. Thus the induced multigraph upstairs is a subgraph of a forest. By F\"urst and Rautenbach \cite[Theorem 2]{FR}, a connected simple subcubic graph with no isolated vertices has an acyclic matching of size at least $n/4$ apart from $K_{2,3}$, $K_4$ with one edge subdivided, and $K_{3,3}$. For $c\ge9$ we have $n\ge7$ so none occurs.

\begin{lemma}\label{colored}
For every connected simple graph of order $n$, cycle rank $e$ and maximum degree at most $\Delta\ge3$, the acyclic matching number is at least $(n-1-(\Delta-2)e)/\Delta$. In particular every connected simple subcubic graph has an acyclic matching of size at least $(n-1-e)/3$.
\end{lemma}
\begin{proof}
Root a spanning tree $T$ and properly color its edges with $\Delta$ colors, processing parents before children. Each endpoint of a non-tree edge has tree degree at most $\Delta-1$. For a non-tree edge whose two endpoints both have tree degree $\Delta-1$, require the colors missing at its ends to differ. Such non-tree edges are vertex-disjoint, since a tree-degree-$(\Delta-1)$ vertex has at most one non-tree incident edge. Set the root's child-edge colors to be distinct. At a later vertex of tree degree $\Delta-1$, its incoming edge is colored. Any of the $\Delta-1\ge2$ other colors can be omitted while coloring its $\Delta-2$ child edges. If its paired endpoint has been processed, choose a missing color different from its partner's. At all other vertices give distinct available colors to the child edges.

For each color $i$ let $M_i$ be its matching in $T$. The endpoints of a non-tree edge can both belong to $V(M_i)$ for at most $\Delta-2$ colors: this holds if one endpoint has lower degree, and otherwise follows from the distinct missing colors. In each affected matching remove an edge incident to one end of the chord, if both ends still occur. The remaining endpoints of each matching induce a subgraph of $T$. Together the matchings retain at least $(n-1)-(\Delta-2)e$ edges. The largest one gives the result.
\end{proof}
Applying Lemma~\ref{colored} to $U$ gives $h\ge(n-1-e)/3\ge(c-3)/3$. Thus $s(L(G))\le c-1-\lceil(c-3)/3\rceil=\lfloor2c/3\rfloor$. The two smaller values $f(2)=1$, $f(3)=2$ follow from the earlier paper and cover the smallest compressions. The existence and lower bound for $\alpha$ were proved in \cite{prev}.

A simple $K$ has a simple cubic expansion, as splitting a vertex while separating incidences never creates a loop or a parallel pair. For $c\ge5$ the expansion has order $J\ge8$. Applying \cite{FR} directly gives $h\ge\lceil J/4\rceil$ and hence $s(L(G))\le\lfloor(c-1)/2\rfloor$. The smaller source bounds are known \cite{prev}. If $U$ is a tree with a matching missing at most one vertex, that matching is acyclic and has size $\lfloor n/2\rfloor$. If it is a cycle, omitting one vertex gives one of size $\lfloor(n-1)/2\rfloor$ and $n\ge c-1$. In either case $h\ge\lfloor(c-2)/2\rfloor$ as required. The alternating polygon chains of \cite{PP} have two end loop vertices and $c-2$ interior double pairs. Their compression is a path at $c\ge4$, and they attain the stated bound. The earlier cases through $c=3$ also cover the smallest degenerate compressions. Finally let $r$ count the deleted loop vertices and contracted double pairs of $C$. The bound $h\ge\lceil(2c-2-r)/4\rceil$ reaches $\lfloor(c-2)/2\rfloor$ whenever $r\le5$ at even $c$, or $r\le7$ at odd $c$.
In fact, let $T$ be any subcubic tree on at least two vertices, and fill each missing cubic incidence by a new loop vertex attached by a single link. The maximum of the free-pin relaxation on this cubic graph is exactly $J-\nu(T)$. The matching bound gives one direction. For the other, pin a minimum cover of $T$ and choose every remaining link and loop odd with positive sign. The unpinned old vertices are independent. Each new loop vertex has diagonal three, and eliminating them leaves diagonal $1-(3-\deg_T(v))/3=\deg_T(v)/3>0$ on each remaining old vertex. Hence all $J-\nu(T)$ surviving directions are positive. This formula concerns the free-pin relaxation.

The coefficient $2/3$ is sharp for the acyclic-matching obstruction itself. Take $m\ge2$ hubs in a path, subdivide each of its $m-1$ links once, and add leaves to bring every hub's tree degree to three. The resulting tree $T$ has $m$ hubs, $m-1$ connectors, and $m+2$ leaves. It has a matching of size $m$ using one private leaf per hub, and the hubs cover every edge. Fill each absent cubic incidence of $T$ by a distinct additional vertex bearing a loop. The resulting cubic kernel has $n=3m+1$ nonloop vertices and $n+2$ loop vertices, so $c=3m+3$. Its qualifying edges are exactly the edges of $T$ and its largest acyclic qualifying matching has size $m$. Thus Lemma~\ref{forest} on this kernel gives exactly $c-1-m=2c/3$. Such multigraphs are realized as branch kernels of simple graphs by subdividing loops. This construction concerns the obstruction's strength; the actual maximum signature on this family remains a separate question.

\section{Exact finite determination}\label{certificate}
Suppose a hypothetical improvement has $\mu\ge J-d$. The computation below uses the integer $d$ as the total allowed loss; the tests below take $d=3$ and $4$. In Lemma~\ref{form}'s first relaxation an independent equation or a payment consumes one unit of the budget. At an intermediate quotient, write $q$ for this combined number. If $A$ is its remaining signed matrix on $m$ coordinates with no further operations, the target is
\[
 m-\sig(A)=2n_-(A)+n_0(A)\le d-q.
\]
In both search procedures, an acyclic qualifying matching of size $d-q+1$ eliminates the entire branch by Lemma~\ref{forest}, even if more operations are performed. The lemma remains applicable with arbitrary nonnegative diagonal weights outside the matched pairs.

\paragraph{Generating the cubic graphs.} A possible counterexample has no acyclic qualifying matching of size $d+1$. Two different procedures generate these exceptional cubic types. For the first, the compression implies $c-2\le n\le4d$, since $n\ge7$ in the cases of interest. Generate the connected simple subcubic $U$ in that range using nauty's \texttt{geng} \cite{nauty}. Discard $U$ containing an acyclic matching of size $d+1$. Restore all possible edge multiplicities subject to degree at most three and $n=c-2+\ell$. Inflate each vertex either to one ordinary cubic vertex or to a doubled pair when its degree permits. Each doubled pair uses different constituent vertices for its two outside slots; permuting them gives isomorphic inflations. Fill every spare slot with a distinct loop vertex. Multiple outside adjacencies cannot have two ordinary endpoints. These choices exhaust the possible preimages, as in \cite{prev}. Discard inflations having an acyclic matching of size $d+1$.

The independent procedure grows all exceptional cubic types by order. Start with the cubic completions of the connected subcubic graphs on four vertices. At each subsequent even order introduce all exceptional simple cubic graphs. Such simple seeds have order at most $4d$, apart from the six-vertex exception of \cite{FR}. For a nonsimple connected cubic graph on at least six vertices, either remove a loop vertex and suppress the resulting degree-two neighbor, or suppress a doubled pair by joining its two outside ends. In both cases a connected cubic multigraph with two fewer vertices results. Conversely the two expansions on one edge $uv$ are (i) subdivide it by $x$ and attach a loop vertex $y$ at $x$, and (ii) replace it by $u x$, two copies of $xy$, and $yv$. The formulas include $u=v$. An acyclic qualifying matching in the smaller graph lifts through either expansion. If the altered edge is outside the matching, its deletion from the induced forest introduces no cycle. If it was matched, replace it by $ux$, dropping $v$ from the selected endpoints. In the second expansion $ux$ is isolated from the remaining selected vertices except through $u$; in the first, its other old neighbor $v$ is unselected. This preserves acyclicity and size. Consequently only exceptional smaller graphs need be expanded.

For order at least four, the underlying simple skeleton of a connected cubic multigraph determines it up to isomorphism. Degree-one vertices get loops and degree-two vertices must be paired with a degree-two neighbor and their common edge doubled. The induced degree-two graph is a disjoint union of paths unless the entire connected skeleton is an even cycle; a path has at most one perfect matching, and the two completions of an even cycle are isomorphic. The second generator therefore canonically labels the ordinary simple skeleton using \texttt{labelg} and removes duplicate strings at each order. The two generators' canonical output files coincide at all orders in the verification table.

\paragraph{Checking a cubic graph.} The first pass uses free coordinate pins from Lemma~\ref{form}. In a signed equation system each connected relation class has one free parameter up to signs or is pinned. Pin every original vertex in a killed class and contract a spanning tree in each surviving one. This costs exactly the original rank, counting every original vertex of a killed class. Redundant equations can be restored as odd edges with a nonnegative term on the class. Free edge signs absorb the choices of coordinate signs. In a class formed by $t$ original vertices the new diagonal starts at $t$; each remaining loop can be chosen positive, adding two. Odd links between classes with multiplicity $k$ give any off-diagonal coefficient in $\{-k,-k+2,\ldots,k\}$. Diagonal switching makes these weights nonnegative on a spanning forest. The program checks all such matrices, for all pin sets and all contraction forests within budget, until the target is eliminated or the graph passes to the tighter check.

The second pass uses only the first relaxation of Lemma~\ref{form}. Try charging any current coordinate; deleting any current nonloop edge; contracting any nonloop edge as an equation; or killing a coordinate by a negative equation on an available loop. Deleting a loop is dominated by keeping it positive. Positive equations on loops are likewise dominated. Every action either costs one or lowers dimension by one. Each pays for its new unit of $q$. These choices include any original order of operations: before a contradictory last equation on a relation cycle one may contract a spanning tree, after which that equation kills the resulting coordinate. All surviving odd signs can be chosen independently after quotienting. Identical labeled matrices of weights and multiplicities at the same cost may be reused. Exhausting all signatures with no successful branch proves the upper bound.

A separate audit reverses each successful relaxed partition. On its pinned original vertices enumerate every internal equation forest. Among the resulting marked classes, choose any that have a remaining loop to be killed by a contradictory equation. Each other marked class must be assigned a charge or deletion incident to it, since such a payment supplied its pin. Enumerate all such assignments, with further incident payments up to the budget, and evaluate the resulting quotient. This covers the preimages of every relaxed partition: assign each original charge to its class and orient each deleted link toward the class chosen to be pinned, and omit payments incident only to genuinely zero classes. This inverse enumeration supplies the tabulated certificate at $c=12$, where it treats all 104,655 successful relaxed partitions. At $c=13$ it independently checks the forward search on 27,034 successful relaxed partitions. Both inverse checks give zero violations.


\paragraph{Exact arithmetic.} The primary inertia test uses symmetric fraction-free elimination. In a trailing matrix with positive denominator $b$, choose a nonzero diagonal pivot $p$ and replace the remaining numerator block by $(|p|A_{ij}-\operatorname{sign}(p)A_{ip}A_{pj})/b$ with new denominator $|p|$ (here $i p$ denotes the pivot column). The divisions are integral by the minor identity for a Schur complement. A zero diagonal block with a nonzero off-diagonal entry receives the unimodular change $e_i\mapsto e_i+e_j$ and so has a nonzero pivot. A zero trailing form gives its whole dimension as nullity. Each negative pivot contributes two loss units. Small integer operations use a checked signed 128-bit intermediate with promotion to GMP for large stored entries. An independent audit computes entire characteristic polynomials by the Faddeev--LeVerrier recurrence and counts positive and negative roots by sign variations. Descartes' bounds are exact for a real-rooted polynomial and its reflection. The test reports its sampling stride and exact counts.

\begin{center}\begin{tabular}{rrrrr}\toprule
$c$ & Allowed loss $d$ & Cubic types & First-pass signed forms & Tighter signed forms\\\midrule
10&3&1,040&228,320&1,070,274\\
11&3&324&63,742&450,829\\
12&4&27,117&171,031,006&4,634,038\\
13&4&10,308&20,409,606&130,896,864\\
14&4&2,401&4,080,741&43,413,192\\
15&4&284&47,528&7,445,764\\
\bottomrule\end{tabular}\end{center}
Each row has zero surviving candidates after the tighter test. The first pass may stop early on a graph which goes to the tighter pass. All counts are matrix evaluations, including repeats. The $d=3$ catalog at $c=12$ contains 52 cubic types and also gives zero final violations. For $d=4$ the two catalogs at $c=16$ are empty. A separate explicit-subset implementation of the forest test agrees on 228,531 target sizes in 45,866 weighted or ordinary graphs. A constructive implementation of Lemma~\ref{colored} was also checked on all 27,505 connected simple subcubic graphs of orders six through twelve.

The upper bound at larger $c$ follows from Theorem~\ref{main} for $d=3$ once $c\ge13$ and for $d=4$ once $c\ge16$. To attain the sharp cases, take a chain of alternating five- and four-cycles starting with a five-cycle; join vertex 1 of one cycle to vertex 0 of the next, labeled consecutively around each cycle. These are the constructions of \cite{PP}. The accompanying verifier builds their actual line graphs and calculates the characteristic polynomials in exact integer arithmetic.

\begin{thebibliography}{9}
\bibitem{PP} A. Paone and M. Paone, \emph{Line-Graph Signature Beyond the 2-Core}, v1.3 (2026), \href{https://doi.org/10.5281/zenodo.21706797}{doi:10.5281/zenodo.21706797}.
\bibitem{FR} M. F\"urst and D. Rautenbach, A lower bound on the acyclic matching number of subcubic graphs, \emph{Discrete Math.} \textbf{341} (2018). \href{https://arxiv.org/abs/1710.10076}{arXiv:1710.10076}, Theorem 2.
\bibitem{prev} Blake Stone's mathematics project, \emph{Line-graph signatures at eight and nine cycles and a uniform upper bound} (2026). \href{https://github.com/blakestone-x/math/blob/5a9d2730dd28ae3031c488960ae2f6cd08a36803/Line-graph-signature/eight-nine.pdf}{Companion paper}, Sections 2--3.
\bibitem{finite} Blake Stone's mathematics project, \emph{A finite kernel for extremal line-graph signature} (2026). \href{https://github.com/blakestone-x/math/blob/5a9d2730dd28ae3031c488960ae2f6cd08a36803/Line-graph-signature/addendum.pdf}{Finite reduction}.
\bibitem{nauty} B. D. McKay and A. Piperno, Practical graph isomorphism, II, \emph{J. Symbolic Comput.} \textbf{60} (2014), 94--112. \href{https://pallini.di.uniroma1.it/}{nauty 2.9.0 and documentation}.
\end{thebibliography}
\end{document}
